\ifdefined\iscombined
	\lecturetitle{Software Quality}
\else
\documentclass{article}
\usepackage{changepage}
\usepackage{centernot}
\usepackage{listings}
\usepackage{amsthm}
\usepackage{braket}
\usepackage{comment}
\usepackage{algorithm}
\usepackage{algpseudocode}
\usepackage{amsmath}
\usepackage{amsfonts}
\usepackage{sectsty}
\usepackage{hyperref}
\usepackage{fancyhdr}
\usepackage[top=1in,bottom=1in,left=1in,right=1in]{geometry}
\usepackage{float}
\usepackage{graphicx}
\usepackage{tabularx}
\usepackage{mathtools}
\usepackage{tikz}
\usetikzlibrary{backgrounds, calc, positioning}

\date{
	\vspace{-.75em}
	{\large \today}
}

\title{
	\vspace*{-1.5em}
	{\Huge Lecture 23} \\
	\vspace{-1em}
	\rule{\textwidth/4}{.01em} \\
	\vspace{-.25em}
	{\Large Software Quality}
	\vspace{-.75em}
}

\author{
	{\large Matthew Farah}
}

\hypersetup{
	colorlinks=true,
	linkcolor=blue,
	filecolor=magenta,
	urlcolor=blue,
	citecolor=blue,
	anchorcolor=blue
}

\fancyhead{}
\fancyhead[L]{\today}
\fancyhead[R]{Lecture 23 - Software Quality}

\setlength{\parindent}{0cm}

\lstset{
	basicstyle=\ttfamily\small,
	frame=single,
	tabsize=4,
	breaklines=true,
	numbers=left,
	numberstyle=\tiny,
	numbersep=8pt,
	xleftmargin=1.5em,
	framexleftmargin=1.5em
}

\newcounter{example}
\renewcommand{\theexample}{\arabic{example}}

\newenvironment{example}[1][]{
	\refstepcounter{example}
	\begin{adjustwidth}{1.5em}{}
	\noindent\textbf{\ifx&#1&Example \theexample:\else Example \theexample \ (#1):\fi}
	\ignorespaces
}{
	\vspace{.5em}
	\end{adjustwidth}
}

\newcounter{subexample}
\renewcommand{\thesubexample}{\alph{subexample}}

\newenvironment{subexample}{
	\setcounter{step}{0}
	\refstepcounter{subexample}
	\vspace{1em}\par\noindent
	\begin{adjustwidth}{1.5em}{}
	\noindent\textbf{(\thesubexample)}
	\ignorespaces
}{
	\par
	\vspace{1em}
	\end{adjustwidth}
}

\newenvironment{solution}{
	\vspace{.5em}
	\par
	\begin{adjustwidth}{1.5em}{}
	\textbf{{Solution:}}
	\par
	\vspace{.5em}
}{
	\vspace{1em}
	\end{adjustwidth}
}

\newcounter{step}
\renewcommand{\thestep}{\Roman{step}}

\newenvironment{step}{
	\refstepcounter{step}
	\vspace{1em}\par\noindent
	\ignorespaces\noindent\textbf{(\thestep)}
	\ignorespaces
}{
	\par
	\vspace{1em}
}

\sectionfont{\fontsize{12}{12}\bfseries}
\subsectionfont{\fontsize{10}{12}\normalfont}
\subsubsectionfont{\fontsize{10}{12}\normalfont}

\begin{document}

\maketitle
\pagestyle{fancy}

\fi

\begin{itemize}
	\item Software quality of a system is a measure of its current and future risk.
	\item The objective of testing is to increase software quality.
	\begin{itemize}
		\item Planning is the process of generating mechanisms to increase software quality.
		\begin{itemize}
			\item Involves creating or using an existing testing strategy.
		\end{itemize}
		\item Monitoring is the process of verifying that such planning effectively increases software quality.
	\end{itemize}
	\item It is better to create imperfect tests early than perfect tests later.
	\begin{itemize}
		\item Acts as a form of review.
		\item Encourages making incremental improvements.
	\end{itemize}
	\item Most defects originate from requirements, not code.
	\begin{itemize}
		\item \emph{Aside: Lowkey cap but whatever ig}.
	\end{itemize}
	%TODO: Ask if anyone has examples of focus risks and scope
	\item All testing strategies have three major components.
	\begin{enumerate}
		\item Focus risks.
		\begin{itemize}
			\item What kinds of risks the testing strategy is attempting to address.
		\end{itemize}
		\item Scope.
		\begin{itemize}
			\item Which business needs the testing strategy is attempting to address.
		\end{itemize}
		\item Details.
		\begin{itemize}
			\item Specifies important considerations of the testing strategy.
			\item For example:
			\begin{itemize}
				\item How risks are identified.
				\item How risks can be minimized.
				\item What metrics can help track risks.
			\end{itemize}
		\end{itemize}
	\end{enumerate}
	\item Different testing strategies optimize software quality with respect to different constraints.
	\item Some examples of widely used testing strategies include:
	\begin{itemize}
		\item The cleanroom strategy.
		\begin{itemize}
			\item Places an emphasis on the identification of risks.
		\end{itemize}
		\item The Software Reliability in Engineering Testing (SERT) strategy.
		\begin{itemize}
			\item Places an emphasis on creating many tests.
		\end{itemize}
		\item The eXtreme Programming (XP) strategy.
		\begin{itemize}
			\item \emph{Aside: Can't remember what she said about this one and it's not on the slides}.
			\item \emph{Aside: According to Google, emphasizes early and incremental automated testing}.
		\end{itemize}
	\end{itemize}
\end{itemize}

\ifdefined\iscombined
	\newpage
\else
	\end{document}
\fi
